% Source: https://artofproblemsolving.com/community/c3h1608948_this_will_take_time
% Stable identity: aops-topic-1608948-q01
Every round removes one contestant, and the game ends when one chump is left. By symmetry, all that matters is how many chumps $c$ and nerds $n$ remain, whether Richard is still in, and whether Richard is the chooser. Let
\[
R(c,n)=P(\text{Richard survives}\mid\text{Richard is the chooser}),\qquad Q(c,n)=P(\text{Richard survives}\mid\text{Richard is in, another chump chooses}),
\]
with $R(1,n)=1$. In a state with $c$ chumps and $n$ nerds the chooser has $c-1+n$ possible opponents.

\textbf{Richard chooses.} He picks a chump with probability $\dfrac{c-1}{c-1+n}$ and must win ($\dfrac{1}{2}$); he picks a nerd with probability $\dfrac{n}{c-1+n}$ and must win ($\dfrac{1}{3}$). After a win he is still the chooser:
\[
R(c,n)=\dfrac{c-1}{c-1+n}\cdot\dfrac{1}{2}\,R(c-1,n)+\dfrac{n}{c-1+n}\cdot\dfrac{1}{3}\,R(c,n-1).
\]

\textbf{Someone else chooses.}
\begin{itemize}
\item The chooser picks Richard (probability $\dfrac{1}{c-1+n}$): Richard must win ($\dfrac{1}{2}$) and becomes the chooser.
\item The chooser picks another chump (probability $\dfrac{c-2}{c-1+n}$): a chump other than Richard leaves and the new chooser is not Richard.
\item The chooser picks a nerd (probability $\dfrac{n}{c-1+n}$): with probability $\dfrac{1}{3}$ nothing changes except that the nerd leaves; with probability $\dfrac{2}{3}$ the chooser leaves too, and the new chooser is Richard with probability $\dfrac{1}{c-1}$.
\end{itemize}
So
\[
Q(c,n)=\dfrac{\dfrac{1}{2}R(c-1,n)+(c-2)\,Q(c-1,n)+n\left[\dfrac{1}{3}Q(c,n-1)+\dfrac{2}{3}\left(\dfrac{R(c-1,n-1)+(c-2)\,Q(c-1,n-1)}{c-1}\right)\right]}{c-1+n},
\]
with $Q(1,n)$ never needed (when $c=2$ the term with $c-2$ vanishes).

\textbf{Computing.} With no nerds, $R(c,0)=\dfrac{1}{2^{c-1}}$, and the recursions give the table
\[
\begin{array}{c|ccccc}
 & c=2 & c=3 & c=4 & c=5 & c=6\\\hline
R(c,0) & \dfrac{1}{2} & \dfrac{1}{4} & \dfrac{1}{8} & \dfrac{1}{16} & \dfrac{1}{32}\\[2mm]
Q(c,0) & \dfrac{1}{2} & \dfrac{3}{8} & \dfrac{7}{24} & \dfrac{15}{64} & \dfrac{31}{160}\\[2mm]
R(c,1) & \dfrac{1}{3} & \dfrac{5}{36} & \dfrac{1}{16} & \dfrac{7}{240} & \dfrac{1}{72}\\[2mm]
Q(c,1) & \dfrac{2}{3} & \dfrac{31}{72} & \dfrac{5}{16} & \dfrac{233}{960} & \dfrac{71}{360}\\[2mm]
R(c,2) & \dfrac{13}{54} & \dfrac{1}{12} & \dfrac{1}{30} & \dfrac{31}{2160} & \dfrac{13}{2016}\\[2mm]
Q(c,2) & \dfrac{41}{54} & \dfrac{11}{24} & \dfrac{29}{90} & \dfrac{2129}{8640} & \dfrac{2003}{10080}
\end{array}
\]
For example, $Q(2,1)=\dfrac{1}{2}\left[\dfrac{1}{2}\cdot1\right]+\dfrac{1}{2}\left[\dfrac{1}{3}\cdot\dfrac{1}{2}+\dfrac{2}{3}\cdot1\right]=\dfrac{1}{4}+\dfrac{5}{12}=\dfrac{2}{3}$: with only Richard and the chooser left, if the chooser loses to the nerd, Richard is the last chump.

The game starts with $c=6$, $n=2$ and Marcus (not Richard) choosing, so
\[
P(\text{Richard survives})=Q(6,2)=\dfrac{2003}{10080}\approx\boxed{0.199}.
\]
